\chapter{Results}
\label{ch:results}

\section{Observed target relations}
\label{sec:membership}

Table~\ref{tab:relations} compares four fixed models on 64 additional
executions of 16 held configurations. REL had the lowest loss and the
highest row accuracy in this comparison, and THM also performed better
than either relation-free model. THM's NLL difference from EVENT GRU
was $-0.149750$ with a Bonferroni-adjusted 97.5\% interval of
$[-0.156307,-0.143693]$, and its difference from flat history attention
was $-0.174032$ with an interval of $[-0.180327,-0.166922]$.
Intervals were computed only for these two THM contrasts, so the
comparison of REL with THM in the table is descriptive. Because the
resampling draws repetitions within the same configurations, the
intervals describe uncertainty over repeated runs under those
configurations.

\begin{table}[tb]
\centering
\caption{Target-aware prediction on 64 new executions and 512 rows from
16 previously evaluated held configurations. NLL is execution-weighted
and accuracy is measured per row. The table averages three seed scores.}
\label{tab:relations}
\begin{tabular}{@{}lrr@{}}
\toprule
Predictor & NLL & Top-1 accuracy \\
\midrule
EVENT GRU & 0.217124 & 89.06\% \\
Flat history attention & 0.241406 & 87.76\% \\
Target history model (THM) & 0.067375 & 97.27\% \\
Relation attention (REL) & \textbf{0.045797} & \textbf{97.98\%} \\
\bottomrule
\end{tabular}
\end{table}

On the separate 64-execution, 704-row evaluation with four targets
and unseen target names, REL scored 0.066618, compared with 0.402682
for EVENT GRU. Training used two- and three-target configurations.
This comparison tests transfer to more targets and new names within
the same task family.

The membership exchange control tests whether the benefit can come
from the group layout alone. THM was trained with either the observed
memberships or exchanged memberships that preserve the number and
sizes of the groups, with each START and its linked END moved
together. On the 16-run development set, correct and exchanged
memberships produced losses of 0.035109 and 0.120010. On the 127-run
reused set, which had been viewed during earlier development, the two
losses were 0.058365 and 0.086855. The corresponding paired 95\%
intervals for correct minus exchanged membership were
$[-0.166172,-0.024583]$ and $[-0.052213,-0.002666]$. Because an
exchanged membership can contradict a visible command path, these
results show the value of the observed membership relative to such
contradictory groupings.

\section{Event-wise history under benign interleaving}

In a separate collection of 111 eligible executions and 1,163 rows,
Local History outperformed the same-capacity Final History control
most consistently when B and C ran concurrently. Table~\ref{tab:local}
reports both training folds. The Bonferroni-adjusted 99.375\%
intervals over eight contrasts for Local minus Final History NLL were
$[-0.364155,-0.152826]$ in the B-fold and $[-0.494571,-0.210046]$ in
the C-fold. With only C running, Final History performed better in the
C-fold, so scores must be reported together with the benign workload
and the training fold.

\begin{table}[tb]
\centering
\caption{Event-wise and final-state history exposure on selected
conditions of the 111-execution collection. B+C has 32 executions and
352 rows, and C-only has 32 executions and 352 rows. Both folds use
the same evaluation executions for a given condition.}
\label{tab:local}
\begin{tabular}{@{}llrrr@{}}
\toprule
Benign workload & Training fold & Local History & Final History & Event-only \\
\midrule
B+C & B & \textbf{0.155520} & 0.409578 & 0.5101 \\
B+C & C & \textbf{0.061857} & 0.409930 & 0.6728 \\
C-only & C & 0.328044 & \textbf{0.136957} & 1.2040 \\
\bottomrule
\end{tabular}
\end{table}
% TODO(사람): Is the "Event-only" column in Table~\ref{tab:local} the EVENT GRU baseline of Table~\ref{tab:relations}, or the event path of the Local History model? If it is EVENT GRU, the column can be renamed. Its values also have four decimals while the other columns have six; please decide the precision.

A further collection of 110 eligible executions compared Local History
with task adaptations of three prior architectures. In the C-fold
B+C condition, the NLLs were 0.071534 for Local History, 0.079912 for
BiTA, 0.193812 for ORTHRUS, and 1.051944 for CL-AP$^2$, and the
adjusted interval for the Local History versus BiTA difference
included zero. In the C-only C-fold condition, the BiTA adaptation
scored 0.125017 against 0.280621 for Local History, with an adjusted
interval for BiTA minus Local History of $[-0.247739,-0.073143]$.
These comparisons concern prediction-stage adaptations that use this
thesis's inputs, labels, and receipt boundary.

\section{Probability averaging and frequency baselines}

A fixed average of a separately trained single-history baseline and
dual-history model reduced NLL from 0.1352 to 0.0802 on flat A+C runs
and from 0.1609 to 0.0778 on nested A+C runs. After retraining with
five new seeds, matched comparisons left the added benefit of the
dual-history architecture uncertain. The observed-history frequency
table improved on the training prior in the original-action cohort,
but in the two A+C layouts all 660 rows had unseen exact keys, so the
table returned the training prior throughout. Appendix~\ref{app:secondary-results}
reports the full comparisons and intervals.
